\chapter{Results}
\label{ch:results}

\section{How to read this chapter}

Every result is labelled with the strength of evidence behind it:
\textbf{Simulated} (measured against a running planning and control stack with
a mock hardware interface), \textbf{Mock-exercised} (the code path ran and its
failure modes were injected, but no physical device was involved), or
\textbf{Hardware} (a physical arm, microcontroller or participant took part).

\textbf{No result in this chapter is of the Hardware tier for the robot.} No
experiment was run with a participant and no logged artefact from a physical
Kinova arm exists. One category of hardware evidence does exist and is
reported in \cref{sec:channels}: recorded master-arm data.

\section{The five-task set}

Nine task families were consolidated to five. The consolidation criterion for
the bimanual tasks was strict: one arm alone must be unable to complete the
task \emph{given unlimited time}. Because the two reachable sets are disjoint
(\cref{ch:workspace}), that can only be achieved in two ways --- a rigid object
spanning the dead band, or two targets requiring simultaneous satisfaction ---
and both are used.

\begin{table}[htbp]
  \centering
  \caption[The five-task set]{The consolidated task set. Every coordinate was
  verified at $N=10$ repeats over the full densified path with both arms
  confirmed at the home pose; total failures zero.}
  \label{tab:tasks}
  \begin{tabular}{llp{62mm}}
    \toprule
    Task & Arms & Why it earns its place \\
    \midrule
    1 Positioning & L, R, both & Fitts characterisation and warm-up; the
      single-versus-both contrast is the attention bottleneck, needing no
      object to exist \\
    2 Pick and place & one at a time & Single-arm manipulation baseline and the
      only discrete success/failure outcome in the set \\
    3 Rigid carry & both & A rigid body held at two points \SI{500}{\milli\metre}
      apart spans the dead band; one arm cannot hold two separated points \\
    4 Compliant carry & both & Identical paths, compliant object; failure is
      nonlinear and abrupt rather than proportional \\
    5 Dual pursuit & both & The two targets lie in disjoint sets, so one arm
      can never reach both; simultaneity is the success criterion \\
    \bottomrule
  \end{tabular}
\end{table}

\todo{MEASURE: every task metric with a participant. The protocols, scenarios,
metrics and success detection exist and have been exercised with scripted
input, but no human has driven the master arm.}

\section{Master-arm channel health and degradation}
\label{sec:channels}

Twelve recordings made on 27 July 2026 constitute the only physical-hardware
data in the project: \num{10890} rows at \SI{100.0}{\hertz} over
\SI{109}{\second}, capturing all fourteen potentiometer channels, two inertial
units, two force-sensitive resistors and two buttons, in six labelled
directions per arm. \emph{(Hardware.)}

Classifying channels between \emph{distinct} sensor updates rather than
between rows --- the distinction matters, as \cref{sec:instrument} explains ---
gives five live channels on the left arm and two on the right at that date.
Re-measuring on 6 August against the same hardware shows substantial
degradation (\cref{fig:degrade}): the left arm's third and fourth channels,
which spanned \SI{234}{\degree} and \SI{175}{\degree} with no dropouts in July,
are classed incoherent in August.

\begin{figure}[htbp]
  \centering
  \includegraphics[width=\textwidth]{figures/channel_degradation.pdf}
  \caption[Progressive channel degradation]{Dropout fraction per channel on two
  dates ten days apart, same hardware. The left arm's shoulder and elbow
  channels degrade from fully live to incoherent.}
  \label{fig:degrade}
\end{figure}

\begin{figure}[htbp]
  \centering
  \includegraphics[width=\textwidth]{figures/channel_health.pdf}
  \caption[Channel health baseline]{The 6 August baseline. Six of fourteen
  channels are usable.}
  \label{fig:channels}
\end{figure}

\subsection{What degraded input costs}

Replaying recorded traces with every coherent channel held at its median
isolates the contribution of the failed channels alone --- a perfectly steady
hand, so every millimetre of output is contributed by the faults. The left
arm's incoherent channels contribute \SI{54.5}{\milli\metre} of commanded tip
motion; the right arm's contribute \SI{0.014}{\milli\metre}.
\emph{(Simulated replay of hardware data.)}

Regressing true reach on the surviving channels shows the left arm retains
\emph{no} reach observable once its shoulder and elbow channels are excluded
($R^{2} = \num{0.133}$, residual \SI{93}{\percent} of the true spread), whereas
the right arm does ($R^{2} = \num{0.986}$). This identifies the two specific
channels whose repair would restore the most capability, and it is the highest
priority hardware action.

\section{Safety layer}

Fourteen distinct faults were injected and all fourteen were handled, meaning
detected, responded to in the designed way, and made visible in the published
state. \emph{(Mock-exercised.)} Faults injected mid-trial into the experiment
runners caused the affected trial to be marked invalid with the cause
recorded, the session to continue, and the analysis to exclude the trial as
pre-registered.

The most consequential defect found was that the dead-man could not detect the
failure that actually occurs, for the reason given in \cref{ch:method}. After
keying it on the source timestamp, the measured trip times were
\SIrange{0.94}{0.99}{\second} for a degraded-but-publishing master and
\SI{0.97}{\second} for a silent one, against a \SI{1.00}{\second} deadline,
with no false positive over \num{243} fresh frames.
\todo{MEASURE: capture these trip traces to file --- the figures above are from
the engineering log and no artefact was retained}

\section{Clearance}

Driving every task and measuring clearance produced the result with the
largest consequence for the programme, developed fully in \cref{ch:workspace}.
Across all recorded runs, \num{48} of \num{87} sit below the
\SI{0.12}{\metre} real-robot floor while passing the planner's collision check
(\cref{fig:clear}).

\begin{figure}[htbp]
  \centering
  \includegraphics[width=0.85\textwidth]{figures/clearance_by_task.pdf}
  \caption[Clearance by task]{Minimum clearance to the wearer across all
  recorded runs. Tasks working at $|x| \ge \SI{0.35}{\metre}$ clear the floor
  comfortably; those in the inboard band do not.}
  \label{fig:clear}
\end{figure}

\section{Recorded verification}

Every task, scenario and control condition was driven in simulation and
recorded from multiple fixed camera angles, with automated checking of object
handling, gripper actuation and wearer clearance. Across all runs no arm link
entered the wearer's collision geometry and no run showed a stationary robot.
\emph{(Simulated.)}

\subsection{A grasp that was not a grasp}

The first round of recordings contained a defect that no automated check
caught and that only viewing the clips revealed: the gripper was never
commanded. The mock hardware interface initialises the driven knuckle at a
fully closed value, and because nothing in the recording path issued a gripper
command, the hand remained shut throughout while the grasp was represented
solely by an attached marker.

The object-attachment check passed, and correctly so: the marker genuinely was
following the gripper. It was following a gripper that never opened. This is
the clearest instance in the project of a verification that was true and
useless.

The recording path now commands the gripper controller directly and executes
the full sequence --- pre-grasp standoff along the approach vector with the
wrist aligned to the object, straight-line approach, closure to the object's
own width rather than to full closure, lift back along the vector --- using the
existing grasp-pose generator, with the marker attaching only once the
\emph{measured} joint indicates closure.

\subsection{Where a top-down grasp is refused}

Exercising the generator produced a quantified refusal that is itself a
result. A top-down grasp requires $|x| \ge \SI{0.30}{\metre}$: zero of nine
candidates solve at \SI{0.25}{\metre} and nine of nine at \SI{0.30}{\metre},
and the failure persists with collision checking disabled, so it is kinematic
infeasibility of the wrist orientation rather than obstruction. Grasps at the
container, at the tray edge and at the tool cradle are therefore refused, and
the recordings show the arm doing what it can with the refusal displayed
rather than pretending otherwise.

\section{The instrument failed more often than the system}
\label{sec:instrument}

A finding that shaped the methodology more than any design decision is that
surprising measurements were more often defects in the measuring instrument
than in the system measured. The instances are numerous enough to be a result
rather than an anecdote.

A noise floor of \num{0.000} on every sensor channel was the recorder
oversampling a \SI{15}{\hertz} sensor at \SI{50}{\hertz}, so
\SIrange{82}{88}{\percent} of adjacent rows were bit-identical. A workspace
statistic varying by a factor of two between identical runs was a summary
statistic sensitive to a rare tail event; on medians the inter-quartile range
is \SI{0.000}{\metre}. A set of verification clips reported as failing were
failing a colour test calibrated against the requested marker colour rather
than the rendered one --- eleven clips whose objects had just been inspected by
eye. A carried object reported as stationary was a centroid dragged towards
fixed overlay text, of which \num{529} of \num{587} matching pixels were
letters. In each case the system was behaving correctly and the instrument was
not.

\section{Unit tests}

Twenty-eight test files are present, including known-answer recovery for every
continuous task metric, the discrete task outcome, the inverse-kinematics
guard convergence and the cross-arm interference fit. Style tests fail and
failed prior to this work, because the codebase uses double quotes throughout;
this is pre-existing and is not a regression.
